% !TeX root = ../../Book2.tex
% 纲要标题：### 18.1 Brauer 等价
% 纲要位置：计划纲要.md，第 1628 行。

\section{Brauer 等价}
\label{b2:sec:1801}

第十六章的矩阵环 Morita 等价给出 \(D\text{-Mod}\simeq M_r(D)\text{-Mod}\)，而第十七章把每个中心单代数写成 \(M_r(D)\)。矩阵大小可以改变代数的维数，却不改变它背后的除代数。固定中心域 \(k\) 及其标量作用，Brauer 等价正是记录这种模范畴的稳定类型；张量积将在这些类上给出可以取逆的运算。本章固定底域 \(k\)，所有中心单代数有限维、非零，代数同构均保持 \(k\) 标量。

% 纲要要点（供目录核对）：
%
% * 中心单代数的 Brauer 等价；以固定底域的线性 Morita 等价解释稳定分类，正式证明非零角代数保持 Brauer 类
% * 张量积给出的群结构；核对加矩阵块的相容性，以包络同构解释反代数真正给出逆元
% * 除代数代表、指数与周期的区分；指数等于最小有限分裂次数，正式证明扩域后指数的两项整除约束
% * 利用反代数解释第一类对合使 Brauer 类的二倍为零，区分周期二与指数二，保留次数、指数及周期各自的意义
%

\subsection{除代数代表与 Brauer 等价}
\label{b2:subsec:1801:01}

若 \(A=M_p(D)\)、\(B=M_q(D)\)，希望略去原来的矩阵重数，就应允许两侧再加不同大小的矩阵块，因为
\[
M_q(A)\cong M_{pq}(D)\cong M_p(B).
\]
要求两侧都加同一个矩阵大小，会保留原来的重数差异，不能表达这种稳定关系。

\begin{definition}\label{b2:def:brauer-equivalence}
设 \(k\) 为域，\(A,B\) 为中心单 \(k\) 代数。若存在正整数 \(r,s\)，使
\[
M_r(A)\cong M_s(B)
\]
为 \(k\) 代数同构，则称 \(A,B\) \term[Brauer-dengjia]{Brauer 等价}[Brauer equivalent]。
\end{definition}

\begin{proposition}\label{b2:prop:brauer-division-representative}
设 \(k\) 为域。每个中心单 \(k\) 代数都同构于某个 \(M_r(D)\)，其中 \(r\geq1\)，\(D\) 为中心是 \(k\) 的有限维除代数。若 \(A\cong M_r(D)\)、\(B\cong M_s(E)\) 为两个中心单 \(k\) 代数，则 \(A,B\) 在 Brauer 意义下等价，当且仅当 \(D\cong E\) 为 \(k\) 代数同构。上述条件也等价于幺左模范畴 \(A\text{-Mod}\) 与 \(B\text{-Mod}\) 之间存在 \(k\) 线性等价，即该等价在每个 Hom 空间上诱导的映射均为 \(k\) 线性。因此 Brauer 等价是等价关系，每个等价类具有唯一到 \(k\) 代数同构的除代数代表。对任意非零幂等元 \(e\in A\)，角代数 \(eAe\) 以 \(e\) 为单位，也是中心单 \(k\) 代数，并与 \(A\) 在 Brauer 意义下等价；其 \(k\) 标量为 \(ce\)、\(c\in k\)。
\end{proposition}
\begin{proof}
中心单代数是有限维单 Artin 环，由\cref{b2:prop:simple-artinian}可写成 \(M_r(D)\)。矩阵环的中心恰为 \(Z(D)\) 中元素组成的标量矩阵，故 \(Z(D)=k\)。这个结构同构保持中心中的标量，故为 \(k\) 代数同构。若 \(A=M_r(D),B=M_s(D)\)，则 \(M_s(A)\cong M_{rs}(D)\cong M_r(B)\)，所以它们 Brauer 等价。

反过来，定义中的矩阵同构给出同一个单环的两种矩阵分解。\cref{b2:thm:artin-wedderburn}及其唯一性说明除环代表同构。该同构可由一个简单模的自同态反环恢复；矩阵同构保持 \(k\) 的中心作用，所以恢复的同构也保持 \(k\)。于是上述充要条件成立，自反、对称和传递均随之得到。

矩阵 Morita 等价由 \(k\) 双模的张量积构造，保持 Hom 空间的 \(k\) 标量，所以同一个 \(D\) 给出上述 \(k\) 线性范畴等价。反过来，等价保持简单对象，并给出其自同态环的 \(k\) 代数同构。两侧各只有一种简单模类型，自同态环分别为 \(D^{\mathrm{op}}\)、\(E^{\mathrm{op}}\)，故 \(D\cong E\) 也保持 \(k\) 标量。

最后，在 \(A=M_r(D)=\operatorname{End}_D(D^r)\) 中，幂等算子 \(e\) 给出右 \(D\) 空间分解 \(D^r=\operatorname{im}e\oplus\ker e\)。取适应该分解的基，将 \(e\) 写成 \(\operatorname{diag}(I_s,0)\)、\(s\ge1\)。限制到像空间给出 \(eAe\cong\operatorname{End}_D(\operatorname{im}e)\cong M_s(D)\)，保持 \(k\) 作用和单位。因此角代数中心单，且仍有除代数代表 \(D\)。
\end{proof}

因此每个 Brauer 类可以用唯一的中心除代数同构类型表示。例如，所有 \(M_n(k)\) 都与 \(k\) 在 Brauer 意义下等价。实四元数除代数 \(\mathbb H\) 却不与 \(\mathbb R\) 等价，因为这两个除代数不同构。Brauer 等价不表示原代数同构：矩阵大小仍然决定实际的维数。

\subsection{张量积与 Brauer 群结构}
\label{b2:subsec:1801:02}

\begin{theorem}\label{b2:thm:brauer-group}
设 \(k\) 为域。中心单 \(k\) 代数的 Brauer 等价类\footnote{每个有限维代数可由某个 \(k^n\) 上的有限乘法表表示，全部这些乘法表组成集合，取同构及稳定等价类后仍为集合。}组成交换群，运算、单位及逆分别为
\[
[A]+[B]=[A\otimes_k B],\qquad 0=[k],\qquad -[A]=[A^{\mathrm{op}}].
\]
这个群称为 \(k\) 的\term[Brauer-qun]{Brauer 群}[Brauer group]，记为 \(\operatorname{Br}(k)\)。
\end{theorem}
\symidx{\(\operatorname{Br}(k)\)：域 \(k\) 的 Brauer 群}
\begin{proof}
要使张量积下降到已经略去矩阵大小的类上，须检查加矩阵块与张量积相容。张量积仍中心单，由\cref{b2:thm:csa-tensor-product}已知。对矩阵单位逐项相乘，得到
\[
M_r(A)\otimes_k M_s(B)\cong M_{rs}(A\otimes_k B).
\]
其映射把 \((E_{ij}a)\otimes(E_{uv}b)\) 送到以 \((i,u),(j,v)\) 为行列指标、条目为 \(a\otimes b\) 的矩阵单位。它保持乘法并把相应的向量空间基一一对应，故为同构。这说明张量运算不依赖稳定代表元的选择。

张量积的结合、交换因子同构与单位同构给出交换群运算所需的结合律、交换律和单位。\cref{b2:thm:csa-enveloping-isomorphism}给出
\[
A\otimes_k A^{\mathrm{op}}\cong\operatorname{End}_k(A),
\]
右边为 \(M_{\dim_kA}(k)\)，代表零类，所以反代数确实通过张量积将 \(A\) 抵消为分裂代数。这正是它给出群逆元的原因。
\end{proof}

若 \(K/k\) 为域扩张，标量扩张保持中心单性，并与张量积相容，所以诱导群同态
\[
\operatorname{Br}(k)\longrightarrow\operatorname{Br}(K),\qquad
[A]\longmapsto[A\otimes_kK].
\]
其核记为 \(\operatorname{Br}(K/k)\)，称为\term[xiangdui-Brauer-qun]{相对 Brauer 群}[relative Brauer group]；其中的类恰由被 \(K\) 分裂的代数代表，因为零类等价于除代数代表为底域，亦即代数本身为全矩阵代数。
\symidx{\(\operatorname{Br}(K/k)\)：由 \(K\) 分裂的 Brauer 类}

\ProseParagraphBreak
任意有限域的 Brauer 群为零。事实上，有限维除代数在有限底域上只有有限多个元素，由\cref{b2:thm:wedderburn-little}为交换域；若中心正是底域，它只能等于底域。代数闭域的 Brauer 群也为零，因为每个中心单代数都已分裂。

\subsection{除代数代表、指数与周期}
\label{b2:subsec:1801:03}

\begin{definition}\label{b2:def:brauer-index-period}
设 \(k\) 为域，\(A\) 为中心单 \(k\) 代数，写 \(A\cong M_r(D)\)，其中 \(r\geq1\)，\(D\) 为中心除 \(k\) 代数。记 \(\dim_kD=d^2\)，定义
\[
\operatorname{ind}(A)=d
\]
并称它为 \(A\) 的\term[zhongxindan-daishu-zhishu]{指数}[index of a central simple algebra]。把 Brauer 类 \([A]\in\operatorname{Br}(k)\) 的加法阶称为 \(A\) 的\term[Brauer-zhouqi]{周期}[period]，记为 \(\operatorname{per}(A)\)；此处按群元素阶的通常约定，暂允许其取值为 \(\infty\)。\cref{b2:thm:period-divides-index}将证明周期总是有限的，且整除指数。这里 \(\deg A=\sqrt{\dim_k A}=rd\)，零 Brauer 类的周期和指数均为一。
\end{definition}
\symidx{\(\operatorname{ind}(A)\)：中心单代数的指数}
\symidx{\(\operatorname{per}(A)\)：中心单代数的周期}

指数和周期都不随矩阵大小改变，次数却随之改变。例如 \(M_3(\mathbb H)\) 的次数为六、指数为二。\eqref{b2:eq:quaternion-conjugation}中的四元数共轭给出 \(\mathbb H\cong\mathbb H^{\mathrm{op}}\)，所以其 Brauer 类的二倍为零；它又不是零类，故周期恰为二。这个例子区分了次数与另两个数，但不表示一般代数的周期都等于指数。

\ProseParagraphBreak
设 \(k\) 为任意域，\(A\) 为中心单 \(k\) 代数。若一个 \(k\) 线性映射 \(\sigma:A\to A\) 满足 \(\sigma(1)=1\)、\(\sigma(xy)=\sigma(y)\sigma(x)\) 及 \(\sigma^2=\operatorname{id}\)，则称它为\term[diyilei-duihe]{第一类对合}[involution of the first kind]。它作为 \(A\to A^{\mathrm{op}}\) 的 \(k\) 代数同构，给出 \([A]=[A^{\mathrm{op}}]=-[A]\)，故 \(2[A]=0\)。这个条件限制的是周期，不能据此断言指数至多为二；\cref{b2:exm:period-two-index-four}给出周期二、指数四的例子。

\begin{proposition}\label{b2:prop:index-splitting-degree}
设 \(k\) 为域，\(A\) 为中心单 \(k\) 代数。若有限域扩张 \(L/k\) 分裂 \(A\)，则 \(\operatorname{ind}(A)\mid[L:k]\)。而且存在次数恰为 \(\operatorname{ind}(A)\) 的有限可分分裂域。因此
\[
\operatorname{ind}(A)
=\min\{\,[L:k]\mid L/k\;\text{为分裂}\;A\;\text{的有限域扩张}\,\}.
\]
\end{proposition}
\begin{proof}
令 \(A=M_r(D)\)、\(\deg D=d\)。\(A_L=M_r(D_L)\) 分裂，当且仅当 \(D_L\) 分裂，这由单环矩阵分解的除代数唯一性得到。此时 \(D\) 通过 \(D_L\cong M_d(L)\) 作用在列空间 \(L^d\) 上。把它看成左 \(D\) 向量空间，设维数为 \(t\)。比较 \(k\) 维数得
\[
d[L:k]=t\dim_kD=td^2,
\]
所以 \(d\mid[L:k]\)。\cref{b2:thm:csa-separable-splitting}还为除代数 \(D\) 构造了一个次数为 \(d\) 的可分子域并证明它分裂 \(D\)，同一个域也分裂 \(A\)，得到存在性。
\end{proof}

这说明指数记录最小可能的有限分裂域次数；周期记录在 Brauer 群中反复张量多少次才能得到零类。一个是分裂扩张的大小，另一个是群运算的阶数，定义上不应混同。


\begin{proposition}\label{b2:prop:index-after-base-change}
设 \(k\) 为域，\(A\) 为中心单 \(k\) 代数，\(L/k\) 为域扩张，令 \(d=\operatorname{ind}(A)\)、\(e=\operatorname{ind}(A_L)\)，其中 \(A_L=A\otimes_kL\)。则 \(e\mid d\)。若 \(L/k\) 有限，则还有
\[
\dfrac de\mid[L:k].
\]
特别地，若 \([L:k]\) 与 \(d\) 互素，则 \(e=d\)。
\end{proposition}
\begin{proof}
以除代数代表 \(D\) 代替 \(A\) 不改变扩域后的 Brauer 类。写 \(D_L\cong M_s(E)\)，其中 \(E\) 为中心除 \(L\) 代数，次数为 \(e\)。比较次数得 \(d=se\)，故 \(e\mid d\)。

若 \(L/k\) 有限，通过含幺嵌入 \(D\hookrightarrow D_L\cong M_s(E)\)，列空间 \(E^s\) 成为左 \(D\) 向量空间。记其 \(D\) 维数为 \(t\)，比较底层 \(k\) 维数，得到
\[
s e^2[L:k]=t d^2=t s^2e^2.
\]
这些是正整数之间的等式，消去 \(se^2\) 得 \([L:k]=ts\)，即 \(d/e=s\) 整除扩域次数。若该次数与 \(d\) 互素，\(s\) 同时整除二者，只能等于一，所以 \(e=d\)。
\end{proof}
